
\section{Introduction}



\looseness=-1
Using classification over a finite vocabulary as the training objective for autoregressive sequence models is an effective approach for naturally discrete modalities such as text, where large-scale Transformer-based~\citep{attentionvaswani} language models such as LLaMa \citep{llama} and GPT-4  \cite{gpt4} have achieved impressive results. To extend this powerful framework to continuous domains such as image, audio, or video, previous work has mostly relied on discretizing signals using lossy compression algorithms~\citep{vqvae}, such that they become akin to text. In particular, neural audio codecs~\citep{soundstream,encodec} have provided discrete representations of audio that are compact enough to allow for high-quality speech~\citep{audiolm,valle} and music \citep{musiclm,musicgen} generation with autoregressive models. In this context, a Residual Vector Quantizer (RVQ)~\citep{soundstream} transforms an audio frame into a coarse-to-fine hierarchy of tokens. As quantization inevitably introduces a perceptual quality loss, generating high-fidelity audio requires increasing the bitrate of audio tokens, which amounts to using deeper hierarchies of RVQ tokens. A consequence of growing the size of the token matrix (along time and token depth) is an additional computational load for the generative model, as the strong dependencies between tokens of the same frame\citep{lemercier} prevent fully parallel generation. The naive approach of flattening the token hierarchy~\citep{audiolm} being prohibitively expensive, ~\citet{musicgen} introduces a delay pattern that conjugates the computational efficiency of parallel generation with a better modeling of inter-token dependencies.~\citet{rq_transformer} and~\citet{uniaudio} furthermore introduce a smaller RQ-Transformer model that is autoregressive along the depth axis, and ~\citet{moshi} combines this approach with the delay pattern. While these methods currently power state-of-the-art generative models for audio~\citep{moshi, hibiki}, the trade-off imposed by residual quantization between quality and computation remains too constraining for generating high-quality audio on edge devices.

\looseness=-1
This motivates an alternative strategy: autoregressive modeling of continuous latents without quantization. Standard variational autoencoders (VAEs) are easier to train, are not affected by issues such as codebook collapse, and can reconstruct audio at higher fidelity for the same latent dimensionality. Pioneering work in the vision domain that autoregressively models continuous sequences includes GIVT \citep{givt} and MAR \citep{autoreg_image_quant}, followed by some larger models in the image domain \citep{fluid, dart} and attempts in the audio domain \citep{salad, sony_trick, ditar, yang2025generativeaudiolanguagemodeling}. In MAR, the authors model the per-token probability distribution with a diffusion model (in the form of a small MLP) conditioned on a latent variable modeled by an autoregressive transformer backbone. 
SALAD \citep{salad} and DiTAR \citep{ditar} adapt MAR-style diffusion heads for text-to-speech modeling, achieving better audio quality than discrete baselines. However, these works are limited to text-to-speech on small-scale and domain-specific datasets, leaving open the question of how well they can adapt to more complex tasks such as speech continuation (without text supervision) and richer audio domains such as music. We apply CALM to 4 tasks which are speech and music continuation as well as text-to-speech and text-to-music.

\looseness=-1
We propose Continuous Audio Language Models (\ours{}) that predict sequences in the latent space of a VAE, bypassing the need for quantization. While we build on the MAR architecture where a transformer backbone uses $(\mathbf{x}^1, \ldots, \mathbf{x}^{s-1})$ to predict an intermediate latent $\mathbf{z}^s$, which then conditions a head (diffusion model) that models $p(\mathbf{x}^s \vert \mathbf{z}^s)$, we find that without further improvements, it fails to generate rich audio content and is slow to sample from. To overcome this, we introduce several key innovations:

\looseness=-1
\textbf{1. Improving quality and stability}: To mitigate error accumulation during inference, we follow \cite{sony_trick} and, during training, inject noise into the long-term context $(\mathbf{x}^1, \ldots, \mathbf{x}^{s-1})$. Additionally, we introduce a short-context transformer that summarizes recent clean latents, providing the sampling head with both coarse long-range context and fine-grained local information. \textbf{2. Diffusion-to-Consistency replacement}: We replace the diffusion model with a continuous consistency model \citep{continuous_consistency} during training, significantly accelerating inference without compromising sample quality. This change reduces the inference time of the sampler head by a factor of up to $\times 20$ in our music experiments and $\times 12$ in our speech experiments compared to the use of a model that uses an RQ-Transformer as head. \textbf{3. Gaussian Temperature sampling}: Temperature control is crucial for high-quality speech generation, yet consistency models lack a formal mechanism for temperature sampling. We present a heuristic that approximates temperature sampling in the consistency setup. \textbf{4. Head batch multiplier}: Sampling multiple noise levels at training time for the same latent highly accelerates training for a small cost.  \textbf{5. Latent Classifier Free Guidance}: we apply Classifier Free Guidance to the latent variable conditioning the consistency head for the conditioned CALM. \textbf{6. Latent Distillation}: once that we have chosen a latent CFG coefficient, we can distill the CFG computation of the backbone into a student backbone while keeping the same sampling head (MLP), hence dividing the batch size by 2 at inference time. As well, we can apply the distillation to a much smaller student backbone. 
%\textbf{5. Classifier Free Guidance}: We introduce two Classifier Free Guidance (CFG) formulas compatible with consistency models. One is applied to the output space and the other to the latent conditioning of the Consistency Head.

Finally, by using these innovations, we introduce \textbf{Pocket TTS} which is a 100M-parameters text-to-speech model that can run faster than real-time on a laptop CPU. We detail the results of \textbf{Pocket TTS} in Sec.\ref{sec:pocket_tts} and in the following technical report: \href{https://kyutai.org/pocket-tts-technical-report}{kyutai.org/pocket-tts-technical-report}.

%\looseness=-1
%Our model generates high-quality speech and music with lower inference time, outperforming state-of-the-art discrete-token autoregressive models in both efficiency and audio fidelity. To our knowledge, this is the first application of autoregressive continuous latent modeling to both domains.


\begin{figure}[t]
\centering
% \begin{minipage}{0.75\textwidth}
    \includegraphics[width=0.8\textwidth]{fig1.pdf}
% \end{minipage}%
% \hfill
% \begin{minipage}{0.24\textwidth}
    \caption{\small Overview of our model. During training, latent vectors $\mathbf{x}^s$ are noised to encourage the backbone Transformer to focus on coarse structure. The consistency head is a consistency model conditioned on the latent variable $\mathbf{z_\text{long}^s}$ produced by the backbone, as well as a short-term context vector $\mathbf{z_\text{short}^s}$ computed from a short-context Transformer applied to the most recent clean latent tokens.}
    \label{fig:fig1}
% \end{minipage}
\end{figure}
